\chapter{Etat de l'art}

\section{Cadre general}

Le clustering d'intentions occupe une place importante dans les systemes de
dialogue orientes tache et, plus largement, dans les applications de service
client. Il s'agit de regrouper automatiquement des enonces selon l'intention
qu'ils expriment, y compris lorsque ces enonces different fortement sur le plan
lexical. Une telle tache presente un interet applicatif evident, puisqu'elle
facilite la construction de taxonomies d'intentions, la structuration de
donnees d'apprentissage et l'analyse de grands volumes de conversations.

La difficulte du probleme tient au fait que la proximite lexicale ne coincide
pas toujours avec la proximite fonctionnelle. Deux formulations proches peuvent
relever d'intentions distinctes, tandis que deux enonces lexicalement eloignes
peuvent correspondre au meme besoin utilisateur. Cette observation limite la
portee des approches reposant exclusivement sur la proximite geometrique entre
representations vectorielles.

Dans ce contexte, l'article \emph{Dial-In LLM: Human-Aligned LLM-in-the-loop
Intent Clustering for Customer Service Dialogues}~\cite{hong2025dialinllm}
propose d'integrer les grands modeles de langage au coeur du processus de
clustering. Afin de situer cette contribution, il convient d'examiner les
travaux anterieurs relatifs au clustering textuel guide par LLM, a la
decouverte d'intentions et a l'evaluation de la coherence semantique.

\section{Clustering textuel guide par les grands modeles de langage}

L'essor des grands modeles de langage a profondement renouvelle les approches
de clustering textuel. Les methodes classiques s'appuient en general sur des
representations vectorielles fixes, puis sur des algorithmes tels que
\emph{K-Means} ou le clustering hierarchique. Des travaux recents montrent
toutefois que les LLM peuvent apporter un signal semantique plus riche, utile
pour guider l'organisation de l'espace de clustering.

Dans cette perspective, plusieurs approches \emph{LLM-guided} ont ete
proposees. \emph{ClusterLLM} montre que des preferences exprimees par un modele
de langage peuvent servir a orienter la structuration des
clusters~\cite{zhang2023clusterllm}. D'autres travaux s'interessent a
l'utilisation des LLM dans des regimes faiblement supervises ou a faible nombre
d'exemples~\cite{viswanathan2024fewshot}. Enfin, certaines methodes
interviennent a un stade de raffinement a posteriori en ciblant les points
frontieres susceptibles de perturber le regroupement, comme dans
\emph{LLMEdgeRefine}~\cite{feng2024llmedgerefine}.

L'ensemble de ces travaux montre que les LLM peuvent contribuer a la
structuration des donnees autrement que comme simples generateurs de texte. Il
reste cependant frequent qu'ils occupent une position peripherique dans le
pipeline, intervenant soit en amont pour enrichir les representations, soit en
aval pour corriger les regroupements~\cite{hong2025dialinllm}. C'est
precisement sur ce point que le papier etudie se distingue.

\section{Decouverte et clustering d'intentions}

Le clustering d'intentions se distingue du clustering textuel general par sa
granularite et par sa finalite. Il ne s'agit pas seulement d'identifier des
themes, mais de retrouver la fonction communicative d'un enonce. Dans les
dialogues orientes tache, cette exigence est determinante, car un meme theme
peut recouvrir plusieurs intentions operationnellement distinctes.

Plusieurs travaux ont explore cette direction. Gung et al. etudient
l'induction d'intentions a partir de conversations dans le cadre des systemes
de dialogue orientes tache~\cite{gung2023intent}. D'autres approches se
placent dans le cadre de la decouverte de nouvelles intentions a partir de
donnees peu ou partiellement annotees. Lin et al. proposent un clustering
adaptatif avec raffinement de clusters~\cite{lin2020discovering}, tandis que
Zhang et al. developpent une strategie de \emph{deep aligned clustering}
destinee a ameliorer la separation semantique entre intentions
emergentes~\cite{zhang2021discovering}. Dans une perspective voisine, Kumar et
al. introduisent une approche semi-supervisee fondee sur un clustering
contrastif profond pour combiner detection et decouverte
d'intentions~\cite{kumar2022intent}.

Des travaux plus recents ont cherche a articuler petits modeles specialises et
LLM pour la decouverte conversationnelle d'intentions
~\cite{liang2024synergizing}. La methode \emph{IDAS}, quant a elle, exploite
des resumes abstraits generes afin d'enrichir la representation semantique des
enonces avant clustering~\cite{deraedt2023idas}.

Ces contributions ont fait progresser la prise en compte de la semantique dans
le clustering d'intentions. Elles laissent cependant ouvertes plusieurs
questions relatives a l'integration explicite d'un jugement de coherence dans
la boucle de clustering et a la robustesse de ces approches sur des corpus de
service client de grande taille~\cite{hong2025dialinllm}.

\section{Limites des approches existantes}

Plusieurs limites recurrentes apparaissent dans cette litterature.

La premiere est la dependance persistante a la proximite entre embeddings. Dans
le cas des intentions conversationnelles, une telle proximite ne garantit ni la
coherence fonctionnelle, ni une bonne separation entre
clusters~\cite{hong2025dialinllm}.

La deuxieme tient au cout computationnel des approches guidees par LLM. Ce cout
peut provenir d'un grand nombre d'appels au modele, de procedures de
raffinement successives ou du recours a des annotations
auxiliaires~\cite{zhang2023clusterllm,deraedt2023idas,viswanathan2024fewshot,feng2024llmedgerefine}.

La troisieme concerne l'evaluation. La \emph{Normalized Mutual Information}
(NMI), bien qu'utile, ne rend pas compte a elle seule de la coherence
semantique des groupes et peut presenter certains biais de
partition~\cite{jerdee2024nmi}. Dans le cadre du clustering d'intentions, la
qualite d'une partition ne se reduit donc pas a son alignement avec une
annotation de reference.

\section{Positionnement de \emph{Dial-In LLM}}

\emph{Dial-In LLM} s'inscrit en reponse directe a ces limites. L'originalite du
papier ne tient pas uniquement a l'usage d'un LLM, mais a la place qui lui est
accordee dans la methode. Le modele de langage n'est pas mobilise comme un
simple outil auxiliaire, mais comme un composant central d'un schema
\emph{LLM-in-the-loop} ou il intervient directement dans l'evaluation de
clusters intermediaires~\cite{hong2025dialinllm}.

Les auteurs introduisent plus precisement deux modules. Le premier est un
evaluateur de coherence semantique, charge de juger si un cluster peut etre
retenu. Le second est un generateur de noms d'intentions, qui produit un
libelle interpretable de type \emph{Action-Objective}. Cette double
architecture fournit a la fois un signal semantique explicite pour la decision
et une couche d'interpretabilite supplementaire.

Cette orientation rejoint des travaux plus anciens sur la coherence semantique
des structures textuelles. Foltz et al. ont montre l'importance de la
coherence dans l'analyse du texte~\cite{foltz1998measurement}, tandis que
Mimno et al. ont souligne son interet dans l'evaluation de structures
thematiques~\cite{mimno2011optimizing}. \emph{Dial-In LLM} prolonge cette idee
en faisant du jugement de coherence un critere operationnel de selection du
nombre de clusters et de validation des resultats
intermediaires~\cite{hong2025dialinllm}.

\section{Apports methodologiques du papier}

Le papier se distingue egalement par plusieurs choix methodologiques precis.

Le premier est le caractere iteratif du clustering. A chaque iteration,
plusieurs valeurs candidates de \texttt{K} sont testees. Les clusters obtenus
sont evalues par le LLM, qui les classe comme \emph{Good} ou \emph{Bad}. Les
groupes juges coherents sont conserves, tandis que les enonces restants sont
reclusterises lors de l'iteration suivante~\cite{hong2025dialinllm}.

Le deuxieme est l'utilisation des noms generes comme support d'une etape de
fusion semantique. Les libelles sont projetes dans un espace hyperspherique,
puis compares au moyen d'une distance geodesique. Les auteurs s'appuient sur
un critere probabiliste inspire des distributions de von Mises-Fisher pour
decider des fusions entre clusters proches
~\cite{fletcher2004principal,banerjee2005clustering,hong2025dialinllm}.

Le troisieme est la prise en compte explicite du contexte conversationnel. Dans
les dialogues de service client, la signification d'un enonce depend du role du
locuteur et de sa position dans l'echange, ce qui rend ce facteur
particulierement important~\cite{hong2025dialinllm}.

\section{Jeux de donnees et protocole d'evaluation}

Le positionnement du papier se renforce enfin par le choix des jeux de donnees
et par son protocole d'evaluation. En plus des benchmarks anglais BANK77,
CLINC(I), MTOP(I) et MASSIVE(I), les auteurs introduisent un corpus chinois de
grande ampleur, construit a partir de plus de 100\,000 appels reels et annote
en 1\,507 clusters d'intentions. Ce corpus est presente comme plus vaste et
plus representatif des usages industriels que les jeux de donnees classiquement
mobilises dans la litterature
~\cite{casanueva2020efficient,larson2019evaluation,fitzgerald2023massive,hong2025dialinllm}.

L'evaluation ne repose pas uniquement sur la NMI. Les auteurs introduisent
egalement un \emph{goodness score}, destine a estimer la proportion de clusters
juges semantiquement coherents. Cette articulation entre mesure externe et
mesure interne constitue l'un des points les plus interessants du protocole
propose~\cite{jerdee2024nmi,hong2025dialinllm}.

\section{Lien avec le present travail}

Cet etat de l'art fournit le cadre conceptuel du projet de reproduction mene
ici. Il permet de situer l'article etudie par rapport aux travaux les plus
directement lies au clustering guide par LLM, a la decouverte d'intentions et a
l'evaluation de la coherence semantique. Cette mise en perspective est
necessaire pour distinguer ce qui releve d'une reproduction du protocole, ce
qui releve d'une adaptation imposee par les ressources disponibles et ce qui
releve des ameliorations proposees dans la suite du rapport.

\section{Conclusion}

L'etat de l'art montre que les travaux recents ont progressivement deplace le
clustering textuel vers des approches davantage sensibles a la semantique,
notamment grace a l'apport des grands modeles de langage. Toutefois, dans la
majorite des cas, ces modeles demeurent situes en marge du processus de
clustering.

\emph{Dial-In LLM} se distingue par une integration beaucoup plus forte du LLM
dans la boucle de decision. Cette caracteristique justifie l'interet
scientifique du present travail de reproduction : l'enjeu n'est pas seulement
de reexecuter un pipeline, mais d'evaluer de maniere critique une proposition
methodologique recente, ses conditions de reproductibilite et la portee de ses
ameliorations possibles.
